\honest{Fake Justification}{Eliezer Yudkowsky}{2007}

Many Christians who have ``stopped really believing,'' I say, still revere the Bible, as a source of ethics or as literature. \nb{The post gives no example or evidence of what these Christians believe.} Sam Harris answers the first defense: it would take five minutes to write a book with a more coherent and compassionate morality. I answer the second: \textsc{The Lord of the Rings} is ``a vastly superior literary work,'' and even \textsc{Harry Potter} is better, as literature and ``as moral philosophy.'' I do not argue for these rankings.

Then a friend buys a \$1 million laptop and says it will stimulate the economy. Buying 5,000 One Laptop Per Child machines would do that and give out laptops as well. This is not motivated stopping, I say, because no search was ever carried out.

As I observed in ``The Bottom Line,'' only the real causes of a decision determine how good it is. The laptop was bought because it was shiny. A justification added later changes nothing unless it is a new search that could change the conclusion. ``Most criticism carried out from a sense of duty,'' I add without evidence, is a token inspection. \nb{A justification of this kind is what the dictionary calls rationalizing; see the notes on ``Rationalization.''}

To justify revering the Bible for its literary quality, you would have to read through candidate books neutrally. That is ``Easy enough if you're not a Christian.'' \nb{The post asks this neutral reading only of Christians. Its own rankings of Tolkien and \textsc{Harry Potter} come with no reading shown.} Of these believers I then say: ``No search ever occurred.'' \nb{The post gives no case. Its setup shows only that the reverence came first, and the next paragraph grants that a later check would count if it could change the conclusion.}

A real check, I grant, does change the real process behind a conclusion. ``But people overestimate, by far, how likely they really are to change their minds.'' \nb{Four weeks earlier, in ``We Change Our Minds Less Often Than We Think,'' Yudkowsky had quoted evidence for this: 23 of 24 colleagues choosing between job offers took the option they had predicted, though their average confidence in the prediction was 66\%.}

Suppose the laptop buyer now says the goal is to give laptops to as many people as possible, and that the shiny laptop is still the best way. I am ``damned suspicious of this amazing coincidence.'' If rolling dice would not likely give the right answer, neither would any other irrational process. ``It's improbable that you used mistaken reasoning, yet made no mistakes.''
